\section{The Topology of Smooth Manifolds}

\subsection{Completeness}

\bdefn

    A Riemannian manifold $M$ is {\emphcolor geodesically complete} iff every maximal geodesic on it is defined over all of $\bR$.

\edefn

Recall that a smooth manifold is metrizable, so we can talk about bounded sets, and also {\emphcolor metrically complete} manifolds which are of
course smooth manifolds which are complete metric spaces.

\blemm

    Let $(M,g)$ be a connected Riemannian manifold such that there is a $p\in M$ for which $\exp_p$ is defined on all of $T_pM$.
    Then for every $q\in M$ there is a minimizing geodesic from $p$ to $q$.

\elemm

\bproof

    Consider the geoedesic sphere $S_\delta(p)$ about $p$ of radius $\delta$.
    This is the image of a compact subset of $T_pM$ under $\exp_p$ and is thus compact.
    Thus $d(\bul,q)$ must attain a minimum on $S_\delta(p)$, let it be at $x_0$.

    Let $v\in T_pM$ of norm $1$ be such that $\exp_p(\delta v)=x$.
    Then define the geodesic $\gamma(t)=\exp(tv)$, and we claim that $\gamma_v$ is a minimizing geodesic from $p$ to $q$.

    We prove this as follows: let $r=d(p,q)$ then we claim that $\gamma(r)=q$.
    Let
    $$ A = \set{s\in[0,r]}[d(\gamma(s),q)=r-s] . $$
    This is nonempty as it contains $0$.
    Furthermore, $A$ is closed as it is the set of points where two continuous maps (to a Hausdorff space) agree.
    We want to show that $r\in A$ since then $d(\gamma(r),q)=0$ and thus $r=q$.

    We now claim that for every $s\in A$ less than $r$, there is an $\epsilon>0$ such that $s+\epsilon\in A$.
    This is sufficient: let $s=\sup A$ and since $A$ is closed necessarily $s\in A$.
    If $s$ were not $r$, then there is some $s+\epsilon\in A$, contradicting $s$'s maximality.

    Now we know that for $s\in A$,
    $$ d(\gamma(s),q) = r - s = d(p,q) - s \implies s = d(\gamma(s),q) - d(p,q) \leq d(p,\gamma(s)) , $$
    where the final inequality is due to the triangle inequality.
    And since $\gamma$ forms a curve of arclength $s_0$ between $p$ and $s_0$, we see that $d(p,\gamma(s))=s$ for all $s\in A$ so $\gamma$
    is minimizing in $A$.

    So let $s_0\in A$, and form a geodesic sphere $S_{\delta'}(\gamma(s_0))$.
    As before let $x_0'\in S_{\delta'}(\gamma(s_0))$ be such that $d(x_0',q)$ is minimal.
    Now,
    $$ d(\gamma(s_0),q) = r - s_0 = d(p,q) - d(p,\gamma(s_0)) . $$
    And
    $$ d(\gamma(s_0),q) = d(\gamma(s_0),x_0') + d(x_0',q) , $$
    as every curve from $\gamma(s_0)$ to $q$ has to interset $S_{\delta'}(\gamma(s_0))$ and $x_0'$ has minimal distance to $q$.
    Thus we see
    $$ d(x_0',q) = r - s_0 - \delta' . $$
    Now we claim that $\gamma(s_0+\delta')=x_0$ and this will complete the proof (of $(*)$) since this implies
    $d(\gamma(s_0+\delta'),q)=r-(s_0+\delta')$ and so $s_0+\delta'\in A$ as desired.

    So we want to show $\gamma(s_0+\delta')=x_0'$.
    Notice that by the triangle inequality,
    $$ d(p,x_0') + d(x_0',q) \geq d(p,q) = r \implies d(p,x_0') \geq r - d(x_0',q) = s_0 + \delta' . $$
    Let $\beta$ be the piecewise smooth curve obtained from $\gamma|_{[0,s_0]}$ and the radial geodesic from $\gamma(s_0)$ to $x_0'$.
    Then we know as $s_0\in A$,
    $$ L_g(\beta) = L_g(\gamma|_{[0,s_0]}) + \delta' = s_0 + \delta' \leq d(p,x_0') . $$
    So there must be equality, and thus it is a minimizing curve, thus a geodesic and by uniqueness $\beta=\gamma|_{[0,s_0+\delta']}$.
    And thus $\gamma(s_0+\delta')=x_0$.
    \qed

\eproof

\bthrm[title={Hopf-Rinow}]

    Let $(M,g)$ be a connected Riemannian manifold.
    Then the following are equivalent:
    \benum
        \item There is a point $p\in M$ for which $\exp_p$ is defined on all of $T_pM$,
        \item (Heine-Borel) if $K\subseteq M$ is closed and bounded (relative to the metric), then it is compact,
        \item $M$ is metrically complete,
        \item $M$ is geodesically complete.
    \eenum

\ethrm

\bproof

    Suppose (1) and $K\subseteq M$ be closed and bounded, so $K\subseteq B_r(p)$ for some $r>0$ in the metric.
    For $x\in K$ we can connect it to $p$ via a minimizing geodesic $\exp_p(tv)$.
    Thus $x=\exp_p(v)$ for $v\in T_pM$ and since this geodesic is minimizing $\norm v=d(x,p)<r$.
    Thus $K$ is contained in $\exp_p(\cl B_r(0))$ and $\cl B_r(0)\subseteq T_pM$ is compact so its image is compact, and since $K$ is closed
    it is compact.
    So we have shown Heine-Borel.

    (2) implies (3) is standard for metric spaces.
    Let $\set{p_n}$ be a Cauchy sequence in $M$, then it is bound so $\cl{\set{p_n}}$ is closed and bound and thus compact by (2).
    Since compact metric spaces are sequentially compact, $p_n$ has a convergent subsequence.
    Cauchy sequences with convergent subsequences converge, so we see that $M$ is metrically complete.

    Now suppose $M$ is metrically complete and let $\gamma$ be a maximal geodesic.
    If it is not defined on all of $\bR$, suppose its domain is bound from above by $s_0\in\bR$.
    Let $s_i$ increase to $s_0$, and let $x_i=\gamma(s_i)$, then $d(x_i,x_j)\leq L_g(\gamma|_{[s_i,s_j]})=\norm{\gamma'}\abs{s_i-s_j}$.
    Since $s_i$ is convergent, it is Cauchy, and thus so is $x_i$.
    By metric completeness, $x_i$ converges to some $p_0\in M$.

    Let $W$ be a uniformly normal neighborhood of $p_0$ for parameter $\delta>0$.
    Take $j$ large enough so that $s_j>s_0-\delta$ and $x_j\in W$.
    So $B_\delta(x_j)$ is a geodesic ball, so let $\sigma(t)=\exp_{x_j}(t\gamma'(s_j))$, i.e.\ it is the radial geodesic with initial
    velocity $\gamma'(s_j)$ starting at $x_j$.
    As it lies in $B_\delta(x_j)$ it is defined for $t\in(-\delta,\delta)$.

    By uniqueness of geodesics, $\gamma(t+s_j)=\sigma(t)$ where defined.
    Thus
    $$ \tilde\gamma(t) = \cases{\gamma(t)&$t<s_0$\cr \sigma(t-s_j)&$t\in(s_j-\delta,s_j+\delta)$} $$
    is smooth, and a geodesic.
    So it extends $\gamma$ to $s_j+\delta>s_0$, contradicting $\gamma$'s maximality.
    Thus $M$ is geodesically complete.
    
    Clearly (4) implies (1).
    \qed

\eproof

\bnote

    We say that a manifold $M$ is {\emphcolor complete} to mean it is geodesically complete.
    If $M$ is connected, then this equivalent to being metrically complete.
    Otherwise this is equivalent to $M$ being metrically complete on each of its connected components, as completeness is equivalent to
    completeness on each of its components.

\enote

\bcoro

    A complete manifold has minimizing geodesics between any two points.

\ecoro

\bcoro

    Compact manifolds are complete.

\ecoro

\bproof

    We can assume the manifold is connected, in which case this follows from Heine-Borel which is trivially satisfied as every closed subset
    of the manifold is compact.
    \qed

\eproof

\subsection{The Hadamard theorem}

\bdefn

    A surjective continuous map between topological spaces $\pi\colon\tilde X\to X$ is a {\emphcolor covering map} if for every $p\in X$ there is a neighborhood $p\in U$
    such that $\pi^{-1}U=\bigsqcup_i\tilde U_i$ for $\tilde U_i\subseteq\tilde X$ open, and $\pi|_{\tilde U_i}\colon\tilde U_i\to U$ is a homeomorphism for all $i$.
    The $\tilde U_i$ are called {\emphcolor sheets}.

    A {\emphcolor smooth covering map} is a covering map between smooth manifolds which is smooth, and the homeomorphisms from the sheets are diffeomorphisms.
    A smooth covering map is a {\emphcolor Riemannian covering map} if it is between Riemannian manifolds and also a local isometry.

\edefn

It can be shown that if $X$ is connected, then the number of sheets over a neighborhood of $p\in X$ is constant for all $p\in X$, call this the {\emphcolor degree} of the covering.
A classic result regarding covering maps is the following.

\bthrm

    If $\pi\colon\tilde X\to X$ is a covering map and $X$ is simply connected (i.e.\ has trivial fundamental group), then $\pi$ has degree $1$.
    Then $\pi$ is a homeomorphism, and a diffeomorphism if $\pi$ is smooth.

\ethrm

\bthrm

    Suppose $(\tilde M,\tilde g)$ and $(M,g)$ are connected Riemannian manifolds with $\tilde M$ complete, and $\pi\colon\tilde M\to M$ is a local isometry.
    Then $M$ is complete and $\pi$ is a Riemannian covering map.

\ethrm

\bproof

    Recall the path-lifting property of covering spaces: for a continuous path $\gamma\colon I\to M$, there is a lift to $\tilde\gamma\colon I\to\tilde M$ such
    that $\pi\circ\tilde\gamma=\gamma$.
    We first show that $\pi$ has the path lifting property for geodesics: for $\pi(\tilde p)=p$ and $\gamma\colon I\to M$ is a geodesic starting at $p$, we can
    lift $\gamma$ uniquely to a geodesic $\tilde\gamma\colon I\to\tilde M$ starting at $\tilde p$.

    Let $v=\gamma'(0)$ and $\tilde v=d\pi_{\tilde p}^{-1}v\in T_{\tilde p}\tilde M$.
    This is well-defined since $\pi$ is a diffeomorphism in a neighborhood of $\tilde p$ and so $d\pi_{\tilde p}$ is an isomorphism.
    Let $\tilde\gamma$ be the geodesic in $\tilde M$ starting at $\tilde p$ with initial velocity $\tilde v$.
    If a lift exists, this must be our lift since
    $$ d\pi_{\tilde p}\circ\tilde\gamma'(0) = \gamma'(0) = v . $$
    Because $\tilde M$ is complete, $\tilde\gamma$ is defined on all of $\bR$, and since $\pi$ is a local isometry $\pi\circ\tilde\gamma$ is a geodesic.
    It has the same starting point and velocity as $\gamma$, so they are equal as desired.

    Now to show that $M$ is complete, let $p$ be a point in $\pi$'s image.
    If $\gamma\colon I\to M$ is any geodesic starting at $p$, then lift $\gamma$ to $\tilde\gamma$, which has an extension to all of $\bR$ as $\tilde M$ is complete.
    Then $\pi\circ\tilde\gamma$ is an extension of $\gamma$ to all of $\bR$.
    By Hopf-Rinow this means $M$ is complete (since $\exp_p$ is defined on all of $T_pM$).

    Now we show $\pi$ is surjective.
    Let $\tilde p\in\tilde M$ and let $p=\pi(\tilde p)$ and let $q\in M$ be arbitrary.
    Since $M$ is complete and connected, there is a unique minimizing unit-speed geodesic $\gamma$ from $p$ to $q$.
    Lift this to a geodesic $\tilde\gamma$ starting at $\tilde p$, and let $r=d_g(p,q)$.
    Then $\pi(\tilde\gamma(r))=\gamma(r)=q$ so $\pi$ is surjective.

    Let $p\in M$ and $U=B_\epsilon(p)$ be the geodesic ball centered at $p$.
    Write $\pi^{-1}p=\set{\tilde p_i}_{i\in A}$, and for each $i$ let $\tilde U_i$ denote the metric ball of radius $\epsilon$ around $\tilde p_i$.
    We claim that $\tilde U_i$ are disjoint.
    For $i\neq j$ let $\tilde\gamma$ be a unit-speed minimizing geodesic from $\tilde p_i$ to $\tilde p_j$.
    Then $\gamma=\pi\circ\tilde\gamma$ is a geodesic starting and ending at $p$ whose length is the same as $\tilde\gamma$'s.
    Since $U$ is a geodesic ball, $\gamma$ must leave and reenter it, as all geodesics lying in $U$ are radial and thus injective.
    So its length must be at least $2\epsilon$, and so $\tilde\gamma$'s length is at least $2\epsilon$ so the distance between $\tilde p_i$ and $\tilde p_j$
    is at least $2\epsilon$, so $\tilde U_i,\tilde U_j$ are disjoint.

    Now we claim $\pi^{-1}U=\bigsqcup_i\tilde U_i$.
    If $\tilde q\in\tilde U_i$ then there is a geodesic of length $<\epsilon$ from $\tilde p_i$ to $\tilde q$, and $\gamma=\pi\circ\tilde\gamma$ is a geodesic of the
    same length from $p$ to $\pi(\tilde q)$, which shows that $\pi(\tilde q)\in U$.
    Thus $\bigsqcup_i\tilde U_i$ is contained in $\pi^{-1}U$.

    Conversely if $q=\pi(\tilde q)\in U$, let $\gamma$ be a minimizing geodesic from $q$ to $p$ which is of length $r<\epsilon$.
    Lift $\gamma$ to $\tilde\gamma$ starting at $\tilde q$, then $\pi(\tilde\gamma(r))=\gamma(r)=p$ so that $\tilde\gamma(r)=\tilde p_i$ for some $i$.
    Then $\tilde q\in\tilde U_i$ as desired.

    Now we claim $\pi|_{\tilde U_i}\colon\tilde U_i\to U$ is a diffeomorphism.
    Since $\pi$ is a local diffeomorphism, clearly this restriction is.
    It is also bijective: for $q\in U$ consider the radial geodesic $p$ to $q$, lift it to $\tilde M$ and map $q$ to its endpoint.
    Since the length of this geodesic is $d_g(p,q)<\epsilon$, this is well-defined and a clear inverse.
    As a bijective local diffeomorphism, the restriction is a diffeomorphism.
    \qed

\eproof

\bcoro

    Suppose $\tilde M,M$ are connected Riemannian manifolds with $\pi\colon\tilde M\to M$ a Riemannian covering map.
    Then $M$ is complete iff $\tilde M$ is.

\ecoro

\bproof

    If $\tilde M$ is complete then by the above theorem so is $M$.
    Conversely if $M$ is complete, let $\tilde p\in\tilde M$ and $\tilde v\in T_p\tilde M$ be arbitrary.
    Let $p=\pi(\tilde p)$ and $v=d\pi_{\tilde p}(\tilde v)$, then the geodesic $\gamma$ starting at $p$ with initial velocity $v$ is defined over all $\bR$.
    Its lift to $\tilde M$ is also a geodesic defined over all $\bR$, and it starts at $\tilde p$ with initial velocity $\tilde v$.
    \qed

\eproof

\bdefn

    Let $\gamma$ be a geodesic segment from $p$ to $q$ in $M$, say $\gamma(a)=p$ and $\gamma(b)=q$.
    We say that $p$ and $q$ are {\emphcolor conjugate} along $\gamma$ if there is a nonzero Jacobi field along $\gamma$ which vanishes at $t=a$ and $t=b$.

\edefn

\bprop

    Let $p\in M$ and $v\in T_pM$ be in the domain of $\exp_p$, and let $q=\exp_p(v)$.
    Then $v$ is a critical point of $\exp_p$ iff $q$ and $p$ are conjugate along $\gamma=\gamma_v$.

\eprop

\bproof

    If $v$ is a critical point of $\exp_p$, then there is a nonzero $w\in T_v(T_pM)$ such that $d(\exp_p)_v(w)=0$.
    Identifying, as usual, $T_v(T_pM)\cong T_pM$, we take $w\in T_pM$.
    Let $\Gamma$ be the variation of $\gamma$ defined by
    $$ \Gamma(s,t) = \exp_p(t(v+sw)) , $$
    and $J(t)=\partial_s\Gamma(0,t)$ be its variation field.
    As $\Gamma$ is a variation through geodesics, $J$ is a Jacobi field.
    And,
    $$ J(1) = \partial_s\Gamma(0,1) = \evpdv s_{s=0}\exp_p(v+sw) = d(\exp_p)_v(w) = 0 . $$
    So $J$ is a non-trivial Jacobi field along $\gamma$ vanishing at $0$ and $1$, so $p$ and $q$ are conjugate.

    Conversely if $p$ and $q$ are conjugate, then let $J$ be a nonzero Jacobi field along $\gamma$ with $J(0)=0$ and $J(1)=0$.
    We showed (see homework) that $J$ is a variation field of a variation of $\gamma$ through geodesics,
    $$ \Gamma(s,t) = \exp_{\sigma(s)}(tV(s)), $$
    for $\sigma(s)$ a smooth curve and $V(s)$ a vector field along $\sigma$, where $\sigma(0)=p$ and $\sigma'(0)=J(0)=0$,
    $V(0)=\gamma'(0)=v$ and $D_sV(0)=D_tJ(0)=w$.
    We can then take $\sigma(s)=p$ and $V(s)=v+sw\in T_pM$ so
    $$ \Gamma(s,t) = \exp_p(t(v+sw)) . $$
    Since $J(t)=\partial_s\Gamma(0,t)$, we have that $w\neq0$ as otherwise $J$ would be zero.
    Then by the same computation as before, $0=J(1)=d(\exp_p)_v(w)$, so $v$ is a critical point of $\exp_p$.
    \qed

\eproof

\blemm

    Suppose $M$ is a Riemannian manifold with nonpositive sectional curvature.
    Then no two points of $M$ are conjugate along any geodesic.

\elemm

\bproof

    Let $\gamma$ be a geodesic segment from $p$ to $q$ in $M$, say $\gamma(0)=p$ and $\gamma(a)=q$.
    Let $J$ be a nonzero Jacobi field along $\gamma$ with $J(0)=0$.
    By the Jacobi equation,
    $$ D_t^2J = -R(J,\dot\gamma)\dot\gamma . $$
    Consider then the sectional curvature in the direction of $J$ and $\dot\gamma$ for any $t$, this is nonpositive by assumption so
    $$ \iprod{D_t^2J,J} = -\iprod{R(J,\dot\gamma)\dot\gamma,J} = -\Rm(J,\dot\gamma,J,\dot\gamma) \geq 0 . $$
    Now consider the smooth map $j\colon [0,a]\to\bR$ given by $j=\iprod{J,J}$.
    Then $j(0)=0$ and
    $$ j' = 2\iprod{D_tJ,J} , j'' = 2\iprod{D_t^2J,J} + 2\iprod{D_tJ,D_tJ} > 0 . $$
    So $j'(0)=0$ and since $j''>0$ we see that $j'$ is increasing and thus always positive.
    Thus $j$ is strictly increasing, which means that since $j(t)>0$ for all $t>0$, in particular $J(t)\neq0$ for all $t>0$.
    \qed

\eproof

\bthrm[title=Hadamard]

    If $(M,g)$ is a complete connected Riemannian manifold with nonpositive sectional curvature, then for every $p\in M$, the map
    $\exp_p\colon T_pM\to M$ is a smooth covering map.
    Thus if $M$ is simply connected it is diffeomorphic to $\bR^n$.

\ethrm

\bproof

    Combining the previous two results, $\exp_p$ has no critical points and is thus an immersion.
    Since the dimension of $T_pM$ and $M$ are the same, this means $\exp_p$ is a local diffeomorphism.
    If we pull back the metric of $M$ to $T_pM$, this gives a metric to $T_pM$ for which $\exp_p$ is a local isometry.

    The straight lines $t\mapsto tv$ in $T_pM$ are geodesics, which means that $T_pM$ is complete.
    This is because for any $w\in T_pM$ and $v\in T_pM\cong T_wT_pM$, the geodesic $tv$ is defined for all $t\in\bR$, so the restricted exponential
    is defined on all of $T_wT_pM$.

    So $T_pM$ is complete, $\exp_p\colon T_pM\to M$ is a local isometry, and thus it is a Riemannian covering map.
    If $M$ is simply connected, then covering maps are diffeomorphisms so $M\cong T_pM\cong\bR^n$.
    \qed

\eproof

\bcoro

    The higher homotopy groups of a complete connected Riemannian manifold are trivial.

\ecoro

\bproof

    As $\exp_p\colon T_pM\to M$ is a covering map, it induces an isomorphism on higher homotopy groups.
    The homotopy groups of $T_pM\cong\bR^n$ are all trivial as $\bR^n$ is contractible.
    \qed

\eproof

\bnote

    A topological space $X$ such that $\pi_n(X)=G$ and $\pi_k(X)=0$ for all $k\neq n$ is called a $K(G,n)$-space, or an {\emphcolor Eilenberg-MacLane space}.
    Eilenberg-MacLane spaces always exist (if $G$ is Abelian for $n>1$), and are unique up to homotopy equivalence.

    We have thus shown that complete connected Riemannian manifolds are Eilenberg-MacLane spaces for $n=1$.

\enote

